Die Systemarchitektur des Health-Monitoring-Systems folgt einer mehrschichtigen Architektur, die aus einem Frontend, einem Backend sowie mehreren externen Datenquellen besteht. Ziel dieser Architektur ist es, Daten aus verschiedenen Plattformen zentral zu sammeln, zu analysieren und anschließend in einem übersichtlichen Dashboard darzustellen.

Durch die klare Trennung der einzelnen Komponenten wird eine modulare und wartbare Systemstruktur erreicht. Jede Schicht des Systems erfüllt eine spezifische Aufgabe und kommuniziert über definierte Schnittstellen mit den anderen Komponenten.


\subsection{Gesamtübersicht der Architektur}
\setauthor{Elias Brandstätter}
Das System besteht aus vier zentralen Ebenen:

\begin{compactitem}
    \item Frontend (Dashboard)
    \item Backend (API und Datenverarbeitung)
    \item Cloud- und Datenebene
    \item Externe Schnittstellen 
\end{compactitem}

Das Frontend dient als Benutzeroberfläche für Mitarbeiter des Unternehmens und ermöglicht die Anzeige sowie Analyse des Status aller Kunden-Apps. Das Backend stellt eine REST-API bereit, über die das Frontend Daten abrufen kann. Gleichzeitig übernimmt das Backend die Verarbeitung und Aggregation der Daten aus verschiedenen Quellen.

Die Cloud- und Datenebene stellt die zentrale Datenbasis dar, während externe Schnittstellen zusätzliche Informationen aus Drittplattformen liefern.


\subsection{Frontend-Schicht}
\setauthor{Zoe Emily Öllinger}
Das Frontend stellt die Benutzeroberfläche für die Mitarbeiter dar. Hier wird die Health-Übersicht und die Detailseite der einzelenen Apps dargestellt.

Das Frontend besteht aus:

\begin{compactitem}
    \item UI \& UX
    \item Health-Übersicht
    \item Detailseite
    \item Kommentar-Sektion
    \item Erkennung von Statusänderungen
\end{compactitem}


\subsection{Backend-Schicht}
\setauthor{Elias Brandstätter}
Das Backend stellt die zentrale Logik des Systems dar und fungiert als Vermittler zwischen Frontend, Datenbank und externen Schnittstellen. Die Backend-Komponente wurde mit FastAPI entwickelt und läuft auf dem ASGI-Server Uvicorn.

Die wichtigsten Aufgaben des Backends sind:

\begin{compactitem}
    \item Bereitstellung einer REST-API für das Frontend
    \item Abrufen und Verarbeiten von Daten aus Google BigQuery
    \item Integration externer APIs
    \item Durchführung von Statusanalysen
    \item Erkennung von Statusänderungen
\end{compactitem}

Das Backend aggregiert Daten aus mehreren Quellen und stellt sie in strukturierter Form bereit. Diese Daten werden anschließend vom Frontend dargestellt.

Ein wichtiger Bestandteil der Backend-Logik ist die Analyse der Statuswerte der einzelnen Apps. Hierbei werden verschiedene Health-Checks durchgeführt und deren Ergebnisse miteinander kombiniert, um einen Gesamtstatus zu bestimmen.

\subsection{Daten- und Cloud-Schicht}
\setauthor{Elias Brandstätter}
Die Datenbasis des Systems liegt in Google BigQuery, wo alle relevanten Monitoring-Daten gespeichert sind - darunter das aggregierte Health-Check-Label der jeweiligen Kunden-Apps, synchronisierte HubSpot-Ticket-Daten sowie App-Konfigurationsinformationen. Das Backend greift auf diese Daten ausschließlich über SQL-Abfragen zu, was eine präzise Datenextraktion ohne unnötigen Overhead ermöglicht.

Ergänzend dazu speichert das Backend täglich einen lokalen Snapshot des aktuellen Systemzustands als CSV-Datei. Diese Snapshots dienen als Vergleichsbasis für die automatische Erkennung von Statusänderungen, ohne dass eine eigene Historiendatenbank betrieben werden muss.

\subsection{Integration externer Systeme}
\setauthor{Elias Brandstätter}
Neben den eigenen Cloud-Ressourcen bindet das System zwei externe Plattformen ein, deren Daten das Bild über den Zustand einer Kunden-App vervollständigen. Asana liefert Informationen über offene Aufgaben, die einem Kunden zugeordnet sind, und ermöglicht das Erstellen neuer Tasks direkt aus dem Dashboard heraus. HubSpot stellt Support-Ticket-Daten bereit, die Aufschluss darüber geben, ob bereits aktive Supportfälle zu einer App existieren. Beide Datenquellen werden im Backend mit den BigQuery-Daten zusammengeführt und in einer einheitlichen Struktur an das Frontend übergeben.

\subsection{Datenfluss im System}
\setauthor{Elias Brandstätter}
Der typische Datenfluss lässt sich anhand des Aufrufs der Health-Check-Übersicht nachvollziehen:

\begin{enumerate}
\item Das Frontend sendet eine HTTP-Anfrage an den \texttt{/health-checks}-Endpunkt, optional ergänzt durch Filterparameter.
\item Das Backend führt an die BigQuery-Tabelle \texttt{health\_monitoring.health\_data} eine SQL-Abfrage aus und lädt parallel die aktuellen HubSpot-Ticketanzahlen sowie offene Asana-Tasks.
\item Die Daten aus allen Quellen werden pro Kunden-Slug zusammengeführt und mit dem zuletzt gespeicherten Snapshot verglichen, um Statusänderungen zu erkennen.
\item Das Backend gibt die aufbereiteten Daten als strukturiertes JSON-Objekt an das Frontend zurück, das die Ergebnisse im Dashboard darstellt.
\end{enumerate}
